\chapter{Martingales II: Closing the Circle}
\begin{center}\itshape
A martingale is a fair game, and a fair game can still cheat you: its value can converge to zero while its average stays at one.\\
Uniform integrability is exactly what forbids the trick. With it come martingales that run backwards in time,\\
de Finetti's theorem on exchangeable sequences, the question of when two infinite experiments can be told apart,\\
and statistical tests that you may check as often as you like.
\end{center}

In March 1939, on the eve of the Second World War, a young Frenchman named Jean Ville defended a thesis with a polemical title,
\emph{\'Etude critique de la notion de collectif}. Its target was Richard von Mises' definition of a random sequence: a ``collective'' was a sequence of
heads and tails whose frequency of heads tends to $\frac12$ along every subsequence picked by a reasonable selection rule. Ville showed that selection rules are
too weak a test of randomness. A gambler is more dangerous: for every property that has probability zero for a fair coin, there is a betting strategy, never
risking more than the gambler owns, whose capital becomes unbounded on every sequence with that property. To describe the gambler's capital Ville borrowed a word from the
casinos, where a \emph{martingale} was the system of doubling the stake after every loss; the dictionary of the Acad\'emie fran\c caise of 1762 already explains
that ``to play the martingale is to always bet all that was lost''. A few months later Ville was mobilized; he spent a year at the front and a year in a
German prison camp. Joseph Doob, who built the general theory in the 1940s, kept Ville's word for his book \emph{Stochastic Processes} of 1953, and it stuck.

Volume~4 proved the convergence theorems for martingales bounded in $L^1$ and in $L^2$, and warned that convergence in mean needs something more. This chapter supplies it, as Volume~4 promised:
uniform integrability, which tells exactly when a martingale is the sequence of predictions of a final outcome. Then we run martingales backwards and obtain the strong law of
large numbers once again, de Finetti's theorem, which explains why a Bayesian may treat exchangeable data as if they came from an unknown coin, Doob's theorem that Bayesian
learning succeeds, and Kakutani's dichotomy on infinite products. We end where Ville began: his inequality for non-negative martingales gives statistical tests that remain valid
however often the experimenter looks at the data, the honest answer to the ``peeking problem'' of online A/B testing.

\begin{reminder}[from Volume 4 and Volume 7]
For an integrable $X$ and a $\sigma$-algebra $\mathcal G$, $\mathbb E[X\mid\mathcal G]$ is the $\mathcal G$-measurable variable with $\mathbb E[\mathbb E[X\mid\mathcal G]\mathbf 1_A]=\mathbb E[X\mathbf 1_A]$ for $A\in\mathcal G$
(Volume~4, Theorem~10.4); it is linear, monotone, satisfies $|\mathbb E[X\mid\mathcal G]|\le\mathbb E[|X|\mid\mathcal G]$ and ``taking out what is known''
(Volume~4, Proposition~10.5). The tower property is Volume~7, Theorem~11.17. A martingale $(M_n)$ for a filtration $(\mathcal F_n)$ is adapted, integrable, with $\mathbb E[M_{n+1}\mid\mathcal F_n]=M_n$ (Volume~4, Definition~10.7); a stopped
martingale is a martingale (Proposition~10.11), and so is a stopped supermartingale a supermartingale. A supermartingale bounded in $L^1$ converges almost surely (Theorem~10.20, through the upcrossing
inequality, Lemma~10.19), and for a non-negative submartingale in $L^2$, $\mathbb E\max_{k\le n}X_k^2\le4\,\mathbb EX_n^2$ (Theorem~10.15).
A family $\mathcal F$ of random variables is \emph{uniformly integrable} if $\sup_{X\in\mathcal F}\mathbb E[|X|\mathbf 1_{\{|X|>M\}}]\to0$ as $M\to\infty$; if $X_n\to X$ in probability, then $\mathbb E|X_n-X|\to0$
if and only if $(X_n)$ is uniformly integrable (Vitali, Volume~7, Theorem~8.16). For an integrable $X$, $\mathbb E[|X|\mathbf 1_A]$ is small when $P(A)$ is small (Volume~7, Lemma~8.15).
Two probabilities that agree on a $\pi$-system agree on the $\sigma$-algebra it generates (Volume~7, Theorem~7.5).
\end{reminder}

\section{Uniformly integrable martingales}

Here is the cheat. Toss a fair coin; at each toss a gambler bets everything on heads. Starting with one euro, after $n$ tosses he has $M_n=2^n$ if all tosses were heads, and
$0$ otherwise. This is a martingale with $\mathbb EM_n=1$ for all $n$, and $M_n\to0$ almost surely: the limit has mean $0$. The mass $2^n$ hides on an event of probability $2^{-n}$,
exactly the escape that uniform integrability forbids.

\begin{lemma}\label{lem:condui}
If $X$ is integrable, the family of all conditional expectations $\mathbb E[X\mid\mathcal G]$, where $\mathcal G$ ranges over the sub-$\sigma$-algebras of $\mathcal F$, is uniformly integrable.
\end{lemma}
\begin{proof}
Let $Y=\mathbb E[X\mid\mathcal G]$ and $A=\{|Y|>M\}$, which belongs to $\mathcal G$. Since $|Y|\le\mathbb E[|X|\mid\mathcal G]$,
\[
\mathbb E[|Y|\mathbf 1_A]\le\mathbb E\bigl[\mathbb E[|X|\mid\mathcal G]\mathbf 1_A\bigr]=\mathbb E[|X|\mathbf 1_A],
\]
and $P(A)\le\frac{\mathbb E|Y|}M\le\frac{\mathbb E|X|}M$ by Markov's inequality. Given $\varepsilon>0$, choose $\delta$ such that $P(A)<\delta$ implies $\mathbb E[|X|\mathbf 1_A]<\varepsilon$ (Volume~7, Lemma~8.15), and then
$M>\mathbb E|X|/\delta$. The bound $\varepsilon$ does not depend on $\mathcal G$.
\end{proof}

\begin{theorem}[uniformly integrable martingales]\label{thm:uimart}
For a martingale $(M_n)$ with respect to $(\mathcal F_n)$, the following are equivalent:
\begin{enumerate}[label=(\alph*)]
\item $(M_n)$ is uniformly integrable;
\item $M_n$ converges almost surely and in $L^1$;
\item there is an integrable $X$ with $M_n=\mathbb E[X\mid\mathcal F_n]$ for all $n$.
\end{enumerate}
In that case the limit $M_\infty$ satisfies $M_n=\mathbb E[M_\infty\mid\mathcal F_n]$ for every $n$.
\end{theorem}
\begin{proof}
(a)$\Rightarrow$(b). A uniformly integrable family is bounded in $L^1$, so $M_n$ converges almost surely to an integrable $M_\infty$ (Volume~4, Theorem~10.20), hence in probability, hence in $L^1$
by Vitali's theorem.

(b)$\Rightarrow$(c). Let $A\in\mathcal F_n$. For $m\ge n$, $\mathbb E[M_m\mathbf 1_A]=\mathbb E[M_n\mathbf 1_A]$ by the martingale property and the tower rule. As $m\to\infty$,
$|\mathbb E[M_m\mathbf 1_A]-\mathbb E[M_\infty\mathbf 1_A]|\le\mathbb E|M_m-M_\infty|\to0$. So $\mathbb E[M_\infty\mathbf 1_A]=\mathbb E[M_n\mathbf 1_A]$ for all $A\in\mathcal F_n$, that is, $M_n=\mathbb E[M_\infty\mid\mathcal F_n]$: take $X=M_\infty$.

(c)$\Rightarrow$(a) is Lemma~\ref{lem:condui}.
\end{proof}

A uniformly integrable martingale is thus a sequence of \emph{predictions}: $M_n$ is the best guess of the final value $M_\infty$ with the information of time $n$. The gambler's
fortune above is not uniformly integrable: it predicts nothing. Which final value does $\mathbb E[X\mid\mathcal F_n]$ converge to? Not $X$ itself, in general, but its best
prediction given everything that will ever be known.

\begin{theorem}[L\'evy's upward theorem, 1937]\label{thm:levyup}
Let $X$ be integrable, $(\mathcal F_n)$ a filtration and $\mathcal F_\infty=\sigma\bigl(\bigcup_n\mathcal F_n\bigr)$. Then
\[
\mathbb E[X\mid\mathcal F_n]\longrightarrow\mathbb E[X\mid\mathcal F_\infty]\qquad\text{almost surely and in }L^1.
\]
\end{theorem}
\begin{proof}
$M_n=\mathbb E[X\mid\mathcal F_n]$ is a martingale by the tower rule, uniformly integrable by Lemma~\ref{lem:condui}, so it converges almost surely and in $L^1$ to some $M_\infty$, and we may take
$M_\infty=\limsup_nM_n$ (where finite, $0$ elsewhere), which is $\mathcal F_\infty$-measurable. Let $\mathcal L$ be the class of $A\in\mathcal F_\infty$ with $\mathbb E[M_\infty\mathbf 1_A]=\mathbb E[X\mathbf 1_A]$. It contains each
$\mathcal F_n$: for $A\in\mathcal F_n$, Theorem~\ref{thm:uimart} gives $\mathbb E[M_\infty\mathbf 1_A]=\mathbb E[M_n\mathbf 1_A]=\mathbb E[X\mathbf 1_A]$. So the two finite signed measures
$A\mapsto\mathbb E[M_\infty\mathbf 1_A]$ and $A\mapsto\mathbb E[X\mathbf 1_A]$ agree on the $\pi$-system $\bigcup_n\mathcal F_n$. Splitting $M_\infty$ and $X$ into positive and negative parts, the measures
$A\mapsto\mathbb E[(M_\infty^++X^-)\mathbf 1_A]$ and $A\mapsto\mathbb E[(X^++M_\infty^-)\mathbf 1_A]$ are finite, agree on that $\pi$-system (which contains $\Omega$), hence agree on $\mathcal F_\infty$ (Volume~7, Theorem~7.5).
So $\mathcal L=\mathcal F_\infty$, and $M_\infty=\mathbb E[X\mid\mathcal F_\infty]$.
\end{proof}

\begin{corollary}[L\'evy's zero--one law]\label{cor:levy01}
If $A\in\mathcal F_\infty$, then $P(A\mid\mathcal F_n)\to\mathbf 1_A$ almost surely.
\end{corollary}
\begin{proof}
Apply Theorem~\ref{thm:levyup} to $X=\mathbf 1_A$, which is $\mathcal F_\infty$-measurable, so that $\mathbb E[X\mid\mathcal F_\infty]=\mathbf 1_A$.
\end{proof}

As information accumulates, every question about the whole future gets answered: the conditional probability of $A$ goes to $0$ or $1$, and it goes to the right one. Kolmogorov's zero--one law
(Theorem~\ref{thm:kolmogorov01}) is a special case: if $X_1,X_2,\dots$ are independent, $\mathcal F_n=\sigma(X_1,\dots,X_n)$ and $A$ is a tail event, then $A$ is independent of $\mathcal F_n$, so
$P(A\mid\mathcal F_n)=P(A)$ for every $n$, and the limit $\mathbf 1_A$ is the constant $P(A)$: so $P(A)\in\{0,1\}$.

\subsection*{Optional stopping without bounds}

Volume~4 proved $\mathbb EM_T=\mathbb EM_0$ for bounded stopping times, for bounded stopped martingales and for steps bounded with $\mathbb ET<\infty$. Uniform integrability gives the
definitive version, with stopping times that may even be infinite. For a stopping time $T$, the \emph{information at time $T$} is the $\sigma$-algebra
\[
\mathcal F_T=\bigl\{A\in\mathcal F_\infty:A\cap\{T=n\}\in\mathcal F_n\text{ for every }n\bigr\}.
\]
It is a $\sigma$-algebra (the conditions are stable under complements, as $A^c\cap\{T=n\}=\{T=n\}\setminus(A\cap\{T=n\})$, and countable unions). If $S\le T$ are stopping times, $\mathcal F_S\subseteq\mathcal F_T$: for $A\in\mathcal F_S$,
\[
A\cap\{T=n\}=\bigcup_{k\le n}\bigl(A\cap\{S=k\}\bigr)\cap\{T=n\}\in\mathcal F_n.
\]
If $M_n\to M_\infty$, we set $M_T=M_\infty$ on $\{T=\infty\}$. Then $M_T$ is $\mathcal F_T$-measurable, since $\{M_T\in B\}\cap\{T=n\}=\{M_n\in B\}\cap\{T=n\}$.

\begin{theorem}[optional stopping for uniformly integrable martingales]\label{thm:uiost}
Let $(M_n)$ be a uniformly integrable martingale with limit $M_\infty$, and $T$ a stopping time with values in $\{0,1,\dots\}\cup\{\infty\}$. Then $M_T$ is integrable and
\[
M_T=\mathbb E[M_\infty\mid\mathcal F_T].
\]
In particular $\mathbb EM_T=\mathbb EM_0$, and $\mathbb E[M_T\mid\mathcal F_S]=M_S$ for stopping times $S\le T$.
\end{theorem}
\begin{proof}
By Theorem~\ref{thm:uimart}, $M_n=\mathbb E[M_\infty\mid\mathcal F_n]$, so $|M_n|\le\mathbb E[|M_\infty|\mid\mathcal F_n]$ and, as $\{T=n\}\in\mathcal F_n$,
\[
\mathbb E|M_T|=\sum_{n\le\infty}\mathbb E\bigl[|M_n|\mathbf 1_{\{T=n\}}\bigr]\le\sum_{n\le\infty}\mathbb E\bigl[|M_\infty|\mathbf 1_{\{T=n\}}\bigr]=\mathbb E|M_\infty|.
\]
Let $A\in\mathcal F_T$. Each $A\cap\{T=n\}$ lies in $\mathcal F_n$, so $\mathbb E[M_n\mathbf 1_{A\cap\{T=n\}}]=\mathbb E[M_\infty\mathbf 1_{A\cap\{T=n\}}]$, also for $n=\infty$. Summing over $n$ (dominated convergence, both
sums being dominated by integrable variables) gives $\mathbb E[M_T\mathbf 1_A]=\mathbb E[M_\infty\mathbf 1_A]$. With the measurability of $M_T$ this is the claim. Taking $A=\Omega$ and comparing with
$T=0$ gives $\mathbb EM_T=\mathbb EM_\infty=\mathbb EM_0$. For $S\le T$, the tower rule and $\mathcal F_S\subseteq\mathcal F_T$ give
$\mathbb E[M_T\mid\mathcal F_S]=\mathbb E[M_\infty\mid\mathcal F_S]=M_S$.
\end{proof}

\begin{example}[a gambler against the house]\label{ex:biasedruin}
A gambler wins $1$ with probability $p$ and loses $1$ with probability $q=1-p$ at each round, with $p<q$: the house has an edge. Starting from $0$ and allowed unlimited credit, what is the
probability that the gambler is ever ahead by $a$, a positive integer? Let $S_n$ be the gain and $T$ the first time $S_n=a$ (possibly $T=\infty$). Then $Z_n=(q/p)^{S_n}$ is a martingale,
since $\mathbb E[(q/p)^{\xi}]=p\cdot\frac qp+q\cdot\frac pq=1$ for a step $\xi$. The stopped martingale $Z_{T\wedge n}$ takes values in $(0,(q/p)^a]$, since $S_{T\wedge n}\le a$, so it is
uniformly integrable. On $\{T=\infty\}$, $S_n\to-\infty$ by the strong law, and $Z_{T\wedge n}\to0$. By Theorem~\ref{thm:uiost} applied to $Z^T$,
\[
1=\mathbb EZ_0=(q/p)^a\,P(T<\infty),\qquad\text{so}\qquad P(T<\infty)=\Bigl(\frac pq\Bigr)^a.
\]
At European roulette, betting on red, $p=\frac{18}{37}$: the chance of ever being $20$ euros ahead is $(18/19)^{20}\approx0.34$.
\end{example}

For the fair game $p=q$, the walk does reach $a$ almost surely, while $\mathbb ES_T=a\neq0=\mathbb ES_0$: the stopped walk is not uniformly integrable (Exercise~\ref{ex:symwalk}).

\section{Martingales that run backwards}

Martingales accumulate information. Sometimes the natural structure \emph{forgets} it: a decreasing sequence of $\sigma$-algebras $\mathcal G_0\supseteq\mathcal G_1\supseteq\cdots$, and
the predictions $\mathbb E[X\mid\mathcal G_n]$ of a fixed variable with less and less information. These \emph{backward martingales} converge even more easily than forward ones: they are
automatically uniformly integrable.

\begin{theorem}[L\'evy's downward theorem]\label{thm:levydown}
Let $X$ be integrable, $\mathcal G_0\supseteq\mathcal G_1\supseteq\cdots$ $\sigma$-algebras and $\mathcal G_\infty=\bigcap_n\mathcal G_n$. Then
\[
\mathbb E[X\mid\mathcal G_n]\longrightarrow\mathbb E[X\mid\mathcal G_\infty]\qquad\text{almost surely and in }L^1.
\]
\end{theorem}
\begin{proof}
Let $M_n=\mathbb E[X\mid\mathcal G_n]$. For fixed $N$, the reversed sequence $M_N,M_{N-1},\dots,M_0$ is a martingale for the increasing filtration $\mathcal G_N\subseteq\cdots\subseteq\mathcal G_0$, by the tower
rule. Let $a<b$. The upcrossings of $[a,b]$ by this reversed sequence are the downcrossings of $[a,b]$ by $M_0,\dots,M_N$; call their number $D_N$. By the upcrossing inequality (Volume~4, Lemma~10.19) applied
to the reversed martingale, $(b-a)\,\mathbb ED_N\le\mathbb E[(M_0-a)^-]\le\mathbb E|X|+|a|$, a bound independent of $N$. By monotone convergence, the total number of downcrossings of $[a,b]$ by $(M_n)$ has finite mean, so it is
finite almost surely. If $\liminf M_n<a<b<\limsup M_n$ there are infinitely many downcrossings; so, taking all rational $a<b$, $M_n$ converges almost surely in $[-\infty,\infty]$. By Fatou's lemma
$\mathbb E\liminf|M_n|\le\mathbb E|X|$, so the limit $M_\infty$ is finite almost surely. The family $(M_n)$ is uniformly integrable by Lemma~\ref{lem:condui}, so $M_n\to M_\infty$ in $L^1$ by Vitali's theorem.

Taking $M_\infty=\limsup_{n\ge m}M_n$ (where finite), $M_\infty$ is $\mathcal G_m$-measurable for every $m$, hence $\mathcal G_\infty$-measurable. For $A\in\mathcal G_\infty$, $A\in\mathcal G_n$ for all $n$, so
$\mathbb E[M_n\mathbf 1_A]=\mathbb E[X\mathbf 1_A]$, and letting $n\to\infty$ in $L^1$ gives $\mathbb E[M_\infty\mathbf 1_A]=\mathbb E[X\mathbf 1_A]$.
\end{proof}

\subsection*{The strong law, by symmetry}

Here is perhaps the most elegant proof of the strong law of large numbers: the empirical mean is a backward martingale. The decreasing information is ``the values of the sequence, but not
the order of the first $n$ of them''.

Let $X=(X_1,X_2,\dots)$ be a sequence of random variables, viewed as a random element of $\R^\N$. For $n\ge1$, let $\mathcal E_n$ be the $\sigma$-algebra of events $\{X\in B\}$ where $B\subseteq\R^\N$ is
measurable and invariant under every permutation of the first $n$ coordinates. (Invariant sets form a $\sigma$-algebra, so $\mathcal E_n$ is one.) These decrease, and $\mathcal E_\infty=\bigcap_n\mathcal E_n$ is
the $\sigma$-algebra of exchangeable events of Chapter~1. (If $A=\{X\in B_n\}$ with $B_n$ invariant under the permutations of the first $n$ coordinates, for each $n$, then $A=\{X\in B\}$ with $B=\bigcup_N\bigcap_{n\ge N}B_n$, which is invariant under every finite permutation.) Any function of $X$ that is symmetric in $X_1,\dots,X_n$, such as $S_n=X_1+\dots+X_n$, or $X_{n+1}$, is $\mathcal E_n$-measurable.
Recall that $X$ is \emph{exchangeable} if its law is invariant under finite permutations: $(X_{\pi(1)},X_{\pi(2)},\dots)$ has the same law as $X$ for every finite permutation $\pi$. An i.i.d.\ sequence is exchangeable.

\begin{lemma}[symmetrization]\label{lem:symm}
Let $X$ be exchangeable, $n\ge k\ge1$, and $h:\R^k\to\R$ measurable with $h(X_1,\dots,X_k)$ integrable. Then
\[
\mathbb E\bigl[h(X_1,\dots,X_k)\mid\mathcal E_n\bigr]=\frac{(n-k)!}{n!}\sum_{(i_1,\dots,i_k)}h(X_{i_1},\dots,X_{i_k}),
\]
the average over all $k$-tuples of distinct indices in $\{1,\dots,n\}$.
\end{lemma}
\begin{proof}
The right side $R$ is symmetric in $X_1,\dots,X_n$, hence $\mathcal E_n$-measurable. Let $A=\{X\in B\}\in\mathcal E_n$ and $(i_1,\dots,i_k)$ distinct indices in $\{1,\dots,n\}$. Choose a permutation $\pi$ of $\{1,\dots,n\}$ with
$\pi(j)=i_j$ for $j\le k$, and let $X^\pi=(X_{\pi(1)},X_{\pi(2)},\dots)$, which has the law of $X$. Since $B$ is invariant, $\{X^\pi\in B\}=\{X\in B\}$, so
\[
\mathbb E\bigl[h(X_{i_1},\dots,X_{i_k})\mathbf 1_A\bigr]=\mathbb E\bigl[h(X^\pi_1,\dots,X^\pi_k)\mathbf 1_{\{X^\pi\in B\}}\bigr]=\mathbb E\bigl[h(X_1,\dots,X_k)\mathbf 1_A\bigr].
\]
Averaging over the $n!/(n-k)!$ tuples gives
\[
\mathbb E[R\mathbf 1_A]=\mathbb E[h(X_1,\dots,X_k)\mathbf 1_A].
\]
\end{proof}

\begin{theorem}[strong law of large numbers, by backward martingales]\label{thm:slllnback}
If $X_1,X_2,\dots$ are i.i.d.\ and integrable, then $\frac{S_n}n\to\mathbb EX_1$ almost surely and in $L^1$.
\end{theorem}
\begin{proof}
By Lemma~\ref{lem:symm} with $k=1$ and $h(x)=x$, $\frac{S_n}n=\mathbb E[X_1\mid\mathcal E_n]$. By Theorem~\ref{thm:levydown}, $\frac{S_n}n$ converges almost surely and in $L^1$ to
$L=\mathbb E[X_1\mid\mathcal E_\infty]$. Every event of $\mathcal E_\infty$ is exchangeable, so it has probability $0$ or $1$ by the Hewitt--Savage law (Theorem~\ref{thm:hewittsavage}); hence $L$ is almost surely
constant (as in Corollary~\ref{cor:tailconst}: $c\mapsto P(L\le c)$ takes only the values $0$ and $1$). The constant is $\mathbb EL=\mathbb EX_1$.
\end{proof}

We proved the strong law in Chapter~2 by Etemadi's method; this proof adds the convergence in $L^1$ for free, and, more importantly, it generalizes at once to averages over pairs, triples and
so on (Exercise~\ref{ex:ustat}): the statistician's $U$-statistics.

\section{De Finetti's theorem}

Bruno de Finetti opened his treatise \emph{Theory of Probability} (1970; English translation 1974) with a provocation set in capital letters:
``PROBABILITY DOES NOT EXIST''. He meant that probability is not a property of coins or dice, like their mass, but a degree of belief of a person; the
abandonment of objective probability, he wrote, was as necessary as the abandonment of ``Phlogiston, the Cosmic Ether, Absolute Space and Time, \ldots, or Fairies and Witches''.
This raises a puzzle. Statisticians speak constantly of ``the unknown probability $\theta$ that this coin falls heads''. If probability does not exist, what are they estimating?

De Finetti's answer, in a memoir of 1931 and in his lectures at the Institut Henri Poincar\'e in 1935 (published in 1937), was a theorem. Suppose your beliefs about a sequence of tosses are
\emph{exchangeable}: the order of the results does not matter to you, only how many heads and tails occur. Then, necessarily, you behave exactly as if there were an unknown parameter $\Theta$,
with some prior law, and the tosses were independent with probability $\Theta$ of heads, given $\Theta$. The ``unknown probability'' is not a physical property; it is the limiting frequency,
and its existence is a theorem about your beliefs.

\begin{theorem}[de Finetti, 1931]\label{thm:definetti}
Let $X_1,X_2,\dots$ be an exchangeable sequence with values in $\{0,1\}$, and $S_n=X_1+\dots+X_n$. Then:
\begin{enumerate}[label=(\alph*)]
\item $\Theta=\lim_n\frac{S_n}n$ exists almost surely;
\item given $\Theta$, the $X_i$ are independent Bernoulli$(\Theta)$: for every $n$ and $x\in\{0,1\}^n$ with $s$ ones,
\[
P\bigl((X_1,\dots,X_n)=x\mid\Theta\bigr)=\Theta^s(1-\Theta)^{n-s};
\]
\item consequently, if $\mu$ is the law of $\Theta$,
\[
P\bigl((X_1,\dots,X_n)=x\bigr)=\int_0^1\theta^s(1-\theta)^{n-s}\,\mu(\dd\theta),
\]
and $\mu$ is the only probability on $[0,1]$ with this property.
\end{enumerate}
\end{theorem}
\begin{proof}
Write $(a)_j=a(a-1)\cdots(a-j+1)$ for the falling factorial. Fix $n$ and $x$ with $s$ ones, and let $m\ge n$. Apply Lemma~\ref{lem:symm} with $k=n$ and $h=\mathbf 1_{\{x\}}$. A tuple of distinct indices
$(i_1,\dots,i_n)$ in $\{1,\dots,m\}$ has $(X_{i_1},\dots,X_{i_n})=x$ when the $s$ positions where $x$ has a one receive distinct indices $i$ with $X_i=1$, in order, and the others distinct indices
with $X_i=0$: there are $(S_m)_s(m-S_m)_{n-s}$ such tuples. So
\begin{equation}\label{eq:urnformula}
P\bigl((X_1,\dots,X_n)=x\mid\mathcal E_m\bigr)=\frac{(S_m)_s\,(m-S_m)_{n-s}}{(m)_n}.
\end{equation}
For $n=1$ and $x=(1)$ this reads $\mathbb E[X_1\mid\mathcal E_m]=\frac{S_m}m$, which by Theorem~\ref{thm:levydown} converges almost surely to $\Theta=\mathbb E[X_1\mid\mathcal E_\infty]$: this is (a). In general, the left side
of~\eqref{eq:urnformula} converges almost surely to $P((X_1,\dots,X_n)=x\mid\mathcal E_\infty)$, again by Theorem~\ref{thm:levydown}. On the right, the numerator is a product of $s$ factors $S_m-j$ and $n-s$ factors $m-S_m-j$, and the
denominator a product of $n$ factors $m-j$, with $j<n$; divided by $m$, each factor tends to $\Theta$, $1-\Theta$ or $1$. So
\[
P\bigl((X_1,\dots,X_n)=x\mid\mathcal E_\infty\bigr)=\Theta^s(1-\Theta)^{n-s}.
\]
Since $\Theta$ is $\mathcal E_\infty$-measurable, conditioning both sides on $\sigma(\Theta)\subseteq\mathcal E_\infty$ (tower rule) gives (b), and taking expectations gives the formula in (c).
For uniqueness: taking $x=(1,\dots,1)$, the formula determines $\int\theta^n\,\mu(\dd\theta)$ for every $n$, hence $\int p\,\dd\mu$ for every polynomial $p$. Polynomials are uniformly dense in the continuous
functions on $[0,1]$ (Weierstrass; Bernstein's proof is Volume~7, Exercise~4.6), so $\int f\,\dd\mu$ is determined for every continuous $f$, and so is $\mu([0,t])$, the decreasing limit of $\int f_k\,\dd\mu$ for
continuous $f_k$ equal to $1$ on $[0,t]$, to $0$ on $[t+\frac1k,1]$ and linear in between. Two probabilities on $[0,1]$ with the same distribution function are equal.
\end{proof}

Formula~\eqref{eq:urnformula} has a concrete meaning: given the first $m$ results as an unordered collection, the first $n$ behave like draws \emph{without replacement} from an urn containing $S_m$ ones and
$m-S_m$ zeros. As $m\to\infty$, drawing without replacement from a huge urn becomes drawing with replacement: independent tosses of a $\Theta$-coin.

\begin{example}[P\'olya's urn again]\label{ex:polyadf}
In Exercise~\ref{ex:polya} we saw that the colours drawn from P\'olya's urn (one red and one green ball to start) are exchangeable, with
\[
P\bigl((X_1,\dots,X_n)=x\bigr)=\frac{s!\,(n-s)!}{(n+1)!}=\int_0^1\theta^s(1-\theta)^{n-s}\,\dd\theta,
\]
the Beta integral. By uniqueness in de Finetti's theorem, $\Theta$ is uniform on $[0,1]$, as Volume~4 (Theorem~10.21) proved directly. The urn, which has a memory and no coin in it at all, cannot be
distinguished by any experiment from the following two-stage procedure: pick $\Theta$ uniformly, then toss a $\Theta$-coin forever. Figure~\ref{fig:polyadf} shows the limits of several runs.
\end{example}

\begin{figure}[ht]
\centering\fig{ch09-polya}
\caption{Twelve runs of P\'olya's urn: the proportion $S_n/n$ of red draws, for $n$ up to $2000$ (logarithmic time axis). Each run settles at its own limit $\Theta$; the limits are uniform on $[0,1]$. De Finetti's
theorem says that, given its limit, each run is a sequence of independent tosses of a $\Theta$-coin.}
\label{fig:polyadf}
\end{figure}

The theorem needs an \emph{infinite} exchangeable sequence. Two draws without replacement from an urn with one red and one green ball are exchangeable, but no mixture of i.i.d.\ coins can make
them always different (Exercise~\ref{ex:finiteexch}). The same proof works for any finite set of values instead of $\{0,1\}$; for real-valued sequences, de Finetti's theorem still holds, with $\Theta$ replaced by a random probability measure, the
limit of the empirical measures (Hewitt and Savage, 1955).

\subsection*{Bayesian learning succeeds}

Take the Bayesian story literally: a parameter $\Theta$ with prior law $\mu$ on $[0,1]$, and given $\Theta$, independent tosses with $P(X_i=1\mid\Theta)=\Theta$. Concretely, let $\Theta,U_1,U_2,\dots$ be independent,
$\Theta$ with law $\mu$ and the $U_i$ uniform on $[0,1]$, and $X_i=\mathbf 1_{\{U_i<\Theta\}}$. After $n$ tosses, the Bayesian's \emph{posterior law} is the conditional law of $\Theta$ given $\mathcal F_n=\sigma(X_1,\dots,X_n)$,
and her estimate is the posterior mean $\mathbb E[\Theta\mid\mathcal F_n]$. Does she learn the truth?

\begin{theorem}[Doob, 1949]\label{thm:doobconsistency}
In this model, $\mathbb E[\Theta\mid\mathcal F_n]\to\Theta$ almost surely and in $L^1$, and for every Borel set $B$, the posterior probability $P(\Theta\in B\mid\mathcal F_n)$ tends to $\mathbf 1_{\{\Theta\in B\}}$ almost surely.
\end{theorem}
\begin{proof}
First, $\frac{S_n}n\to\Theta$ almost surely. Indeed, by Fubini's theorem (Volume~7, Theorem~9.6) applied to the independent pair $\Theta$ and $(U_i)$,
\[
P\Bigl(\frac1n\sum_{i\le n}\mathbf 1_{\{U_i<\Theta\}}\to\Theta\Bigr)=\int_0^1P\Bigl(\frac1n\sum_{i\le n}\mathbf 1_{\{U_i<\theta\}}\to\theta\Bigr)\mu(\dd\theta)=1,
\]
by the strong law for each fixed $\theta$. So $\Theta$ is almost surely equal to $\limsup_nS_n/n$, an $\mathcal F_\infty$-measurable variable, and $\mathbb E[\Theta\mid\mathcal F_\infty]=\Theta$. Theorem~\ref{thm:levyup} gives
the first claim, and Corollary~\ref{cor:levy01} applied to $\{\Theta\in B\}\in\mathcal F_\infty$ (up to a null set) gives the second.
\end{proof}

The theorem is a statement about \emph{almost every} $\theta$ for the prior $\mu$, and the caveat is real: David Freedman showed in 1963 that for parameters outside a set of prior measure one,
Bayes estimates can fail to be consistent, once the parameter space is infinite-dimensional. For a coin, all is well: if $\mu$ has a positive density, the posterior concentrates around every $\theta$.
Consistency can be put to work in \emph{Thompson sampling}, proposed by William R.\ Thompson in 1933 for allocating patients to treatments and now used to choose which advertisement to
show: sample $\Theta$ from each posterior, play the best-looking arm, update.

\begin{think}
A sequence is exchangeable and you observe $X_1=\dots=X_{10}=1$. Is the probability that $X_{11}=1$ greater or smaller than $P(X_1=1)$? Show that $\operatorname{Cov}(X_1,X_2)=\operatorname{Var}\Theta\ge0$:
in an infinite exchangeable sequence, outcomes can never be negatively correlated. Laplace's rule of succession, $P(X_{n+1}=1\mid S_n=s)=\frac{s+1}{n+2}$, is the case of a uniform prior.
\end{think}

\section{Kakutani's dichotomy}

You are given an infinite sequence of tosses and told that it comes either from a fair coin or from a coin whose bias at toss $n$ is $\frac12+\varepsilon_n$, with $\varepsilon_n\to0$. With the whole infinite sequence in
hand, can you decide without error? If $\varepsilon_n=\frac1{10}$, certainly: the frequency tells. If $\varepsilon_n=2^{-n}$, certainly not: the two coins differ only in the first few tosses, and those are not enough. Shizuo Kakutani
proved in 1948 that there is nothing in between: two infinite product measures are either perfectly distinguishable or not distinguishable even with probability one.

Two probabilities $P$ and $Q$ on the same space are \emph{equivalent} if they have the same null sets, and \emph{mutually singular}, $P\perp Q$, if there is an event $A$ with $P(A)=0$ and $Q(A)=1$ (Volume~7,
Chapter~10). For probabilities $p$ and $q$ on $\R$ with $q\ll p$, the \emph{Hellinger affinity} is
\[
\rho(p,q)=\int\sqrt{\frac{\dd q}{\dd p}}\;\dd p\in[0,1],
\]
where $\frac{\dd q}{\dd p}$ is the Radon--Nikodym derivative (Volume~7, Theorem~11.4). By Cauchy--Schwarz, $\rho\le\bigl(\int\frac{\dd q}{\dd p}\dd p\bigr)^{1/2}=1$, with equality only if $\frac{\dd q}{\dd p}=1$, that is $p=q$. If $p$ and $q$ are equivalent, $\rho>0$
and $\rho(p,q)=\rho(q,p)$, since both equal $\int\sqrt{fg}\,\dd\lambda$ for densities $f$, $g$ of $p$, $q$ with respect to $\lambda=p+q$.

\begin{theorem}[Kakutani, 1948]\label{thm:kakutani}
For each $n$, let $p_n$ and $q_n$ be equivalent probabilities on $\R$, and let $P=\bigotimes_np_n$ and $Q=\bigotimes_nq_n$ be the product probabilities on $\R^\N$ (Volume~7, Theorem~12.8). Then $P$ and $Q$ are
either equivalent or mutually singular:
\[
\prod_{n=1}^\infty\rho(p_n,q_n)>0\ \Longrightarrow\ P\text{ and }Q\text{ are equivalent};\qquad\prod_{n=1}^\infty\rho(p_n,q_n)=0\ \Longrightarrow\ P\perp Q.
\]
In the first case, $\frac{\dd Q}{\dd P}=\prod_{n=1}^\infty f_n(X_n)$, an almost surely convergent product, where $f_n=\frac{\dd q_n}{\dd p_n}$ and $X_n$ is the $n$-th coordinate.
\end{theorem}
\begin{proof}
Work under $P$, where the coordinates $X_n$ are independent with laws $p_n$, with $\mathcal F_n=\sigma(X_1,\dots,X_n)$. Let $\rho_n=\rho(p_n,q_n)$ and
\[
M_n=\prod_{k\le n}f_k(X_k).
\]
Since $\mathbb E_Pf_k(X_k)=q_k(\R)=1$, $M_n$ is a non-negative martingale with mean $1$. For $B\subseteq\R^n$ a product of Borel sets, $Q((X_1,\dots,X_n)\in B)=\mathbb E_P[M_n\mathbf 1_{\{(X_1,\dots,X_n)\in B\}}]$, by
independence; both sides are probabilities in $B$, so by uniqueness (Volume~7, Theorem~7.5) this holds for all Borel $B\subseteq\R^n$:
\begin{equation}\label{eq:densityn}
Q(A)=\mathbb E_P[M_n\mathbf 1_A]\qquad\text{for }A\in\mathcal F_n.
\end{equation}

\emph{Case $\prod\rho_n>0$.} The variables $N_n=\prod_{k\le n}\sqrt{f_k(X_k)}/\rho_k$ form a non-negative martingale, with $\mathbb E_PN_n^2=\prod_{k\le n}\rho_k^{-2}\le c^{-2}$, where
$c=\prod_{k=1}^\infty\rho_k>0$ (the partial products decrease to $c$). By Doob's inequality (Volume~4, Theorem~10.15) and monotone convergence, $\mathbb E_P\sup_nN_n^2\le4c^{-2}$. Since
$M_n=N_n^2\prod_{k\le n}\rho_k^2\le N_n^2$, the martingale $(M_n)$ is dominated by an integrable variable, so it is uniformly integrable and converges almost surely and in $L^1$ to some $M_\infty$, with
$M_n=\mathbb E_P[M_\infty\mid\mathcal F_n]$ (Theorem~\ref{thm:uimart}). The probability $Q'(A)=\mathbb E_P[M_\infty\mathbf 1_A]$ agrees with $Q$ on every $\mathcal F_n$, by~\eqref{eq:densityn}, hence on $\sigma(\bigcup_n\mathcal F_n)$, the product
$\sigma$-algebra. So $Q\ll P$ with density $M_\infty$. Exchanging the roles of $p_n$ and $q_n$, which does not change $\rho_n$, gives $P\ll Q$.

\emph{Case $\prod\rho_n=0$.} Now $\mathbb E_P\sqrt{M_n}=\prod_{k\le n}\rho_k\to0$. Choose $n_1<n_2<\cdots$ with $\mathbb E_P\sqrt{M_{n_j}}\le4^{-j}$, and let $A_j=\{M_{n_j}>1\}\in\mathcal F_{n_j}$. Then
$P(A_j)\le\mathbb E_P\sqrt{M_{n_j}}\le4^{-j}$ by Markov's inequality applied to $\sqrt{M_{n_j}}$, and by~\eqref{eq:densityn},
\[
Q(A_j^c)=\mathbb E_P\bigl[M_{n_j}\mathbf 1_{\{M_{n_j}\le1\}}\bigr]\le\mathbb E_P\sqrt{M_{n_j}}\le4^{-j},
\]
since $m\le\sqrt m$ for $0\le m\le1$. Let $A=\limsup_jA_j$. By Borel--Cantelli, $P(A)=0$. And $A^c=\bigcup_i\bigcap_{j\ge i}A_j^c$, where $Q(\bigcap_{j\ge i}A_j^c)\le\inf_{j\ge i}4^{-j}=0$; so $Q(A^c)=0$, and $Q(A)=1$.
\end{proof}

\begin{example}[shifted Gaussians and biased coins]\label{ex:kakutaniex}
(a) Let $p_n=\mathcal N(0,1)$ and $q_n=\mathcal N(a_n,1)$. Completing the square, $\varphi(x)\varphi(x-a)=\frac1{2\pi}e^{-(x-a/2)^2-a^2/4}$, so
\[
\rho(p_n,q_n)=\int\sqrt{\varphi(x)\varphi(x-a_n)}\,\dd x=e^{-a_n^2/8}.
\]
So the shifted sequence $(Z_n+a_n)$ of independent standard Gaussians has a law equivalent to that of $(Z_n)$ if $\sum a_n^2<\infty$, and singular otherwise. This is the discrete shadow of the Cameron--Martin
theorem for Brownian motion: shifting a Brownian path by a function $h$ gives an equivalent law exactly when $h'$ is square integrable.

(b) Let $p_n$ be the fair coin and $q_n$ the coin with $P(1)=\frac12+\varepsilon_n$. Then
\[
\rho(p_n,q_n)=\frac{\sqrt{1+2\varepsilon_n}+\sqrt{1-2\varepsilon_n}}2=1-\frac{\varepsilon_n^2}2+O(\varepsilon_n^4).
\]
Since $\prod(1-u_n)>0$ if and only if $\sum u_n<\infty$ (for $0\le u_n<1$), the biased sequence can be detected with certainty if and only if $\sum\varepsilon_n^2=\infty$. For $\varepsilon_n=n^{-1/2}$ it can, although the bias
tends to $0$; for $\varepsilon_n=n^{-0.51}$ it cannot.
\end{example}

\section{Peeking at the data: Ville's inequality}

An online shop tests a new design of its checkout page. Half the visitors see the old page, half the new one, and the analyst computes, every morning, the usual test statistic: with $n$ visitors per
group and $D_n$ the difference in the number of purchases, she declares the new page better as soon as the standardized difference exceeds $1.96$, a test ``at level $5\%$''. Each single test does have level $5\%$.
But she looks every day, and stops at the first success. By the law of the iterated logarithm (Theorem~\ref{thm:hartmanwintner}), even when the two pages are identical, the standardized difference exceeds
$1.96$, indeed exceeds any fixed level, infinitely often, almost surely. Given enough patience, the analyst is \emph{certain} to declare a winner that does not exist. Numerical simulation measures how fast:
with daily looks over a long experiment, the false alarm rate is several times the advertised $5\%$.

Ville's thesis contains the remedy. A non-negative supermartingale cannot often become large.

\begin{theorem}[Ville's inequality, 1939]\label{thm:ville}
Let $(M_n)_{n\ge0}$ be a non-negative supermartingale and $c>0$. Then
\[
P\Bigl(\sup_{n\ge0}M_n\ge c\Bigr)\le\frac{\mathbb EM_0}c.
\]
\end{theorem}
\begin{proof}
Let $T$ be the first $n$ with $M_n\ge c$ ($T=\infty$ if there is none), a stopping time. The stopped process $M_{T\wedge n}$ is a supermartingale (Volume~4, Proposition~10.11), so
\[
c\,P(T\le n)\le\mathbb E\bigl[M_{T\wedge n}\mathbf 1_{\{T\le n\}}\bigr]\le\mathbb EM_{T\wedge n}\le\mathbb EM_0,
\]
using $M_{T\wedge n}=M_T\ge c$ on $\{T\le n\}$ and $M\ge0$. Let $n\to\infty$: $P(T<\infty)\le\mathbb EM_0/c$.
\end{proof}

Think of $M_n$ as the capital of a gambler who bets against the hypothesis, starting with one euro, at fair odds if the hypothesis is true. Ville's inequality says that the chance that the
gambler ever multiplies his euro by $20$ is at most $5\%$. So ``reject the hypothesis as soon as the gambler's capital reaches $1/\alpha$'' is a test of level $\alpha$, \emph{whatever the stopping rule}: one may look
after every observation, stop when one likes, or never stop. Here is such a gambler for the fair coin.

\begin{theorem}[an anytime-valid test of a coin]\label{thm:anytime}
Let $X_1,X_2,\dots$ be tosses with values in $\{0,1\}$, $S_n=X_1+\dots+X_n$, and
\[
M_n=\int_0^1\frac{\theta^{S_n}(1-\theta)^{n-S_n}}{2^{-n}}\,\dd\theta=\frac{2^n\,S_n!\,(n-S_n)!}{(n+1)!}.
\]
\begin{enumerate}[label=(\alph*)]
\item If the coin is fair, $(M_n)$ is a martingale with $M_0=1$, so $P(M_n\ge1/\alpha\text{ for some }n)\le\alpha$ for every $\alpha\in(0,1)$.
\item If the tosses are i.i.d.\ with $P(X_i=1)=\theta\neq\frac12$, then $M_n\to\infty$ almost surely: the test rejects the fair coin eventually, with probability $1$.
\end{enumerate}
\end{theorem}
\begin{proof}
(a) The second expression is the Beta integral. From it, $M_{n+1}/M_n$ equals $\frac{2(S_n+1)}{n+2}$ if $X_{n+1}=1$ and $\frac{2(n-S_n+1)}{n+2}$ if $X_{n+1}=0$; for a fair toss independent of $\mathcal F_n$,
\[
\mathbb E\bigl[M_{n+1}/M_n\,\big|\,\mathcal F_n\bigr]=\frac12\cdot\frac{2(S_n+1)+2(n-S_n+1)}{n+2}=1.
\]
Apply Theorem~\ref{thm:ville} with $c=1/\alpha$.

(b) For $0\le k\le n$, $\binom nkt^k(1-t)^{n-k}\le1$ at $t=\frac kn$ gives $\binom nk\le e^{nH(k/n)}$, where $H(t)=-t\log t-(1-t)\log(1-t)$ is the entropy in nats. So
\[
M_n=\frac{2^n}{(n+1)\binom n{S_n}}\ge\frac{\exp\bigl(n(\log2-H(S_n/n))\bigr)}{n+1}.
\]
By the strong law $S_n/n\to\theta$, and $H(\theta)<\log2$ for $\theta\neq\frac12$ (Gibbs' inequality, Volume~4, Proposition~4.16: $\log2-H(t)$ is the relative entropy of the $t$-coin from the fair coin). So the numerator grows exponentially, the denominator only linearly,
and $M_n\to\infty$.
\end{proof}

\begin{figure}[ht]
\centering\fig{ch09-ville}
\caption{The capital $M_n$ of the betting test of Theorem~\ref{thm:anytime}, logarithmic scale, for $n\le3000$. Grey: ten runs with a fair coin; Ville's inequality says that at most $5\%$ of such runs ever reach the
rust line $1/\alpha=20$. Charcoal: ten runs with $\theta=0.58$; each crosses the line and goes on growing exponentially.}
\label{fig:ville}
\end{figure}

The betting test pays a modest price for its freedom: with a fixed sample size, the classical test needs fewer observations to detect a given bias. But the analyst who cannot resist looking, or who must
stop early when the evidence is overwhelming, has no honest alternative. Such \emph{e-values}, \emph{confidence sequences} and \emph{anytime-valid} tests, going back to Ville and to Herbert Robbins'
work of 1967--1970, are today a lively part of statistics, driven in part by online experiments (the survey of Ramdas, Gr\"unwald, Vovk and Shafer is a good entry point). Exercise~\ref{ex:confseq} turns the test into a confidence sequence: an interval for an unknown
bias that is valid at all times simultaneously.

\begin{think}
In Theorem~\ref{thm:anytime}, the prior on $\theta$ was uniform. What happens with a prior concentrated near $\frac12$? Near $0.9$? A gambler who knows the true $\theta$ would bet with the likelihood ratio
$\prod(2\theta)^{X_i}(2-2\theta)^{1-X_i}$: how fast does its capital grow, and why is this growth rate the exponent of the Chernoff--Stein lemma (Theorem~7.9)?
\end{think}

\section{Exercises}

\begin{exercise}[$\star$]\label{ex:doubling}
Let $\xi_1,\xi_2,\dots$ be independent fair bits and $M_n=\prod_{k\le n}2\xi_k$. Check directly that $(M_n)$ is not uniformly integrable, and that $M_n\neq\mathbb E[M_\infty\mid\mathcal F_n]$.
\end{exercise}

\begin{exercise}[$\star$]\label{ex:symwalk}
Let $S_n$ be the simple symmetric random walk started at $0$ and $T$ the first time it reaches $1$. Using recurrence (Volume~4), $T<\infty$ almost surely. Show that $(S_{T\wedge n})$ is not uniformly integrable,
and deduce from case~(c) of optional stopping (Volume~4, Theorem~10.12) that $\mathbb ET=\infty$.
\end{exercise}

\begin{exercise}[$\star$]\label{ex:finiteexch}
(a) Let $(X_1,X_2)$ take the values $(1,0)$ and $(0,1)$ with probability $\frac12$ each. Show that it is exchangeable but not of the form of Theorem~\ref{thm:definetti}~(c).
(b) Show that for an infinite exchangeable binary sequence, $\operatorname{Cov}(X_1,X_2)=\operatorname{Var}\Theta\ge0$.
\end{exercise}

\begin{exercise}[$\star\star$]\label{ex:polyabeta}
P\'olya's urn starts with $a$ red and $b$ green balls ($a,b\ge1$ integers); each drawn ball is returned with one more of its colour. Show that the sequence of colours is exchangeable and that the limiting proportion of
red has the Beta$(a,b)$ law, with density $\frac{(a+b-1)!}{(a-1)!\,(b-1)!}\theta^{a-1}(1-\theta)^{b-1}$.
\end{exercise}

\begin{exercise}[$\star\star$]\label{ex:ustat}
Let $X_1,X_2,\dots$ be i.i.d.\ and $h:\R^2\to\R$ symmetric with $h(X_1,X_2)$ integrable. Show that
\[
U_n=\binom n2^{-1}\sum_{1\le i<j\le n}h(X_i,X_j)\longrightarrow\mathbb Eh(X_1,X_2)
\]
almost surely and in $L^1$. Example: for $h(x,y)=\frac12(x-y)^2$, $U_n$ is the usual unbiased sample variance.
\end{exercise}

\begin{exercise}[$\star\star$]\label{ex:kakutanipoisson}
(a) Show that the Hellinger affinity of the Poisson laws with means $\lambda$ and $\mu$ is $e^{-(\sqrt\lambda-\sqrt\mu)^2/2}$. When are the laws of two sequences of independent Poisson variables, with means
$\lambda_n$ and $\mu_n$, equivalent?
(b) In Example~\ref{ex:kakutaniex}~(a) with $\sum a_n^2<\infty$, show that $\frac{\dd Q}{\dd P}=\exp\bigl(\sum_na_nX_n-\frac12\sum_na_n^2\bigr)$, the series converging almost surely under $P$.
\end{exercise}

\begin{exercise}[$\star\star$]\label{ex:confseq}
For $\theta_0\in(0,1)$ let
\[
M_n(\theta_0)=\frac{\int_0^1\theta^{S_n}(1-\theta)^{n-S_n}\,\dd\theta}{\theta_0^{S_n}(1-\theta_0)^{n-S_n}},
\]
and let $C_n=\{\theta_0:M_n(\theta_0)<1/\alpha\}$. Show that if the tosses are i.i.d.\ with bias $\theta$, then $P(\theta\in C_n\text{ for all }n)\ge1-\alpha$. Show that $C_n$ is an interval.
\end{exercise}

\section{Hints}

\begin{hint}{ex:doubling}
$\mathbb E[M_n\mathbf 1_{\{M_n>K\}}]=1$ as soon as $2^n>K$. What is $M_\infty$?
\end{hint}
\begin{hint}{ex:symwalk}
If the family were uniformly integrable, Theorem~\ref{thm:uiost} would give $\mathbb ES_T=0$. If $\mathbb ET$ were finite, Volume~4's case~(c) would give the same.
\end{hint}
\begin{hint}{ex:finiteexch}
(a) A mixture would need $\int\theta^2\dd\mu=0$ and $\int(1-\theta)^2\dd\mu=0$. (b) Use Theorem~\ref{thm:definetti}~(c) with $n=2$.
\end{hint}
\begin{hint}{ex:polyabeta}
Compute the probability of a sequence with $s$ reds: its numerators are $a(a+1)\cdots(a+s-1)$ and $b\cdots(b+n-s-1)$ in some order. Compare with the Beta integral.
\end{hint}
\begin{hint}{ex:ustat}
Lemma~\ref{lem:symm} with $k=2$ gives $U_n=\mathbb E[h(X_1,X_2)\mid\mathcal E_n]$.
\end{hint}
\begin{hint}{ex:kakutanipoisson}
(a) Sum $e^{-(\lambda+\mu)/2}\frac{(\sqrt{\lambda\mu})^k}{k!}$. (b) $M_n=\exp(\sum_{k\le n}a_kX_k-\frac12\sum_{k\le n}a_k^2)$; the series $\sum a_kX_k$ converges almost surely, for example by Volume~4, Theorem~10.18.
\end{hint}
\begin{hint}{ex:confseq}
For the true $\theta$, $M_n(\theta)$ is a martingale with mean $1$. For the interval, look at $\log M_n(\theta_0)$ as a function of $\theta_0$.
\end{hint}

\section{Solutions}

\begin{solution}{ex:doubling}
\idea All the mean sits on a single event of probability $2^{-n}$, which vanishes in the limit.
\writeup $M_n=2^n$ on the event that $\xi_1=\dots=\xi_n=1$, of probability $2^{-n}$, and $0$ elsewhere. For any $K$ and $n$ with $2^n>K$, $\mathbb E[M_n\mathbf 1_{\{M_n>K\}}]=2^n\cdot2^{-n}=1$, so
$\sup_n\mathbb E[M_n\mathbf 1_{\{M_n>K\}}]=1$ for every $K$: not uniformly integrable. Almost surely some $\xi_k=0$, so $M_\infty=0$, and $\mathbb E[M_\infty\mid\mathcal F_n]=0\neq M_n$ on an event of probability $2^{-n}>0$.
\end{solution}

\begin{solution}{ex:symwalk}
\idea The walk reaches $1$ for sure, so the stopped value has mean $1$, not $0$: the fair game was beaten, and the only way this is possible is that the stopped martingale is not uniformly integrable and the waiting time has infinite mean.
\writeup $S_T=1$ almost surely, so $\mathbb ES_T=1\neq0=\mathbb ES_0$. If $(S_{T\wedge n})$ were uniformly integrable, Theorem~\ref{thm:uiost} applied to the martingale $S^T$ (whose limit is $S_T$) would give
$\mathbb ES_T=0$. If $\mathbb ET<\infty$, Volume~4's case~(c), with steps bounded by $1$, would also give $\mathbb ES_T=0$. Both are impossible. (The losses before $T$ hide in rare excursions far below $0$.)
\end{solution}

\begin{solution}{ex:finiteexch}
\idea Two draws from a two-ball urn are perfectly negatively correlated; mixtures of coins only produce non-negative correlation.
\writeup (a) The law gives probability $\frac12$ to $(1,0)$ and to $(0,1)$, a set invariant under the swap, so it is exchangeable. If $P(X_1=X_2=1)=\int\theta^2\dd\mu$ and $P(X_1=X_2=0)=\int(1-\theta)^2\dd\mu$ for some $\mu$,
both would be $0$, forcing $\mu$ to be carried by $\{\theta:\theta=0\}\cap\{\theta:\theta=1\}=\emptyset$. (b) By Theorem~\ref{thm:definetti}~(c), $P(X_1=1)=\mathbb E\Theta$ and $P(X_1=X_2=1)=\mathbb E\Theta^2$, so
$\operatorname{Cov}(X_1,X_2)=\mathbb E\Theta^2-(\mathbb E\Theta)^2=\operatorname{Var}\Theta$.
\end{solution}

\begin{solution}{ex:polyabeta}
\idea As in Exercise~\ref{ex:polya}, the probability of a sequence depends only on its number of reds; then recognize a Beta integral and use uniqueness in de Finetti's theorem.
\writeup When $r$ reds and $g$ greens have been drawn, the urn holds $a+r$ reds and $b+g$ greens. So a sequence of $n$ draws with $s$ reds has probability
\[
\frac{a(a+1)\cdots(a+s-1)\cdot b(b+1)\cdots(b+n-s-1)}{(a+b)(a+b+1)\cdots(a+b+n-1)},
\]
which depends only on $s$: the sequence is exchangeable. With $B(x,y)=\int_0^1\theta^{x-1}(1-\theta)^{y-1}\dd\theta=\frac{(x-1)!\,(y-1)!}{(x+y-1)!}$ for positive integers (Volume~5, Section~7.4 gives $B(x,y)$ in general), the probability equals
\[
\frac{(a+s-1)!\,(b+n-s-1)!\,(a+b-1)!}{(a-1)!\,(b-1)!\,(a+b+n-1)!}=\frac{B(a+s,b+n-s)}{B(a,b)}=\int_0^1\theta^s(1-\theta)^{n-s}\,\beta_{a,b}(\theta)\,\dd\theta,
\]
with $\beta_{a,b}$ the Beta$(a,b)$ density. By uniqueness in Theorem~\ref{thm:definetti}, $\Theta=\lim S_n/n$ has the Beta$(a,b)$ law; and the proportion of red balls in the urn, $\frac{a+S_n}{a+b+n}$, has the same limit.
\end{solution}

\begin{solution}{ex:ustat}
\idea Averaging over pairs is the best prediction of one pair given the unordered sample; as the sample grows, the backward martingale converges, and Hewitt--Savage makes the limit constant.
\writeup By Lemma~\ref{lem:symm} with $k=2$, $\mathbb E[h(X_1,X_2)\mid\mathcal E_n]$ is the average of $h(X_i,X_j)$ over ordered pairs of distinct indices, which is $U_n$ by symmetry of $h$. By Theorem~\ref{thm:levydown}, $U_n$
converges almost surely and in $L^1$ to $\mathbb E[h(X_1,X_2)\mid\mathcal E_\infty]$, which is almost surely constant by Theorem~\ref{thm:hewittsavage}, equal to its mean $\mathbb Eh(X_1,X_2)$. For
$h=\frac12(x-y)^2$, $\sum_{i<j}(X_i-X_j)^2=n\sum_i(X_i-\bar X)^2$, so $U_n=\frac1{n-1}\sum_i(X_i-\bar X)^2$, and the limit is $\operatorname{Var}X_1$ when $X_1\in L^2$.
\end{solution}

\begin{solution}{ex:kakutanipoisson}
\idea Affinities of exponential families are explicit; the density of the product is the limit of the martingale of partial likelihood ratios.
\writeup (a) $\rho=\sum_k\sqrt{e^{-\lambda}\frac{\lambda^k}{k!}e^{-\mu}\frac{\mu^k}{k!}}=e^{-(\lambda+\mu)/2}\sum_k\frac{(\sqrt{\lambda\mu})^k}{k!}=e^{-(\lambda+\mu)/2+\sqrt{\lambda\mu}}=e^{-(\sqrt\lambda-\sqrt\mu)^2/2}$. For positive means the laws are equivalent,
and by Theorem~\ref{thm:kakutani} the products are equivalent if and only if $\sum_n(\sqrt{\lambda_n}-\sqrt{\mu_n})^2<\infty$, and singular otherwise.
(b) Here $f_n(x)=\varphi(x-a_n)/\varphi(x)=e^{a_nx-a_n^2/2}$, so $M_n=\exp(\sum_{k\le n}a_kX_k-\frac12\sum_{k\le n}a_k^2)$. Under $P$, $\sum_{k\le n}a_kX_k$ is a martingale with
$\mathbb E(\sum_{k\le n}a_kX_k)^2=\sum_{k\le n}a_k^2$, bounded, so it converges almost surely (Volume~4, Theorem~10.18). Hence $M_n\to\exp(\sum a_nX_n-\frac12\sum a_n^2)$, and Theorem~\ref{thm:kakutani} identifies this limit as
$\frac{\dd Q}{\dd P}$.
\end{solution}

\begin{solution}{ex:confseq}
\idea For the true parameter, the betting capital is a non-negative martingale with mean $1$; Ville's inequality bounds the chance that it ever reaches $1/\alpha$, that is, the chance that the truth ever leaves the set.
\writeup Under bias $\theta$, $M_{n+1}(\theta)/M_n(\theta)$ equals $\frac{S_n+1}{(n+2)\theta}$ if $X_{n+1}=1$ and $\frac{n-S_n+1}{(n+2)(1-\theta)}$ otherwise (as in Theorem~\ref{thm:anytime}), with conditional mean
\[
\theta\cdot\frac{S_n+1}{(n+2)\theta}+(1-\theta)\cdot\frac{n-S_n+1}{(n+2)(1-\theta)}=1.
\]
So $M_n(\theta)$ is a non-negative martingale with $M_0(\theta)=1$, and $P(\theta\notin C_n\text{ for some }n)=P(\sup_nM_n(\theta)\ge1/\alpha)\le\alpha$ by Theorem~\ref{thm:ville}. Finally,
$\log M_n(\theta_0)=\text{const}-S_n\log\theta_0-(n-S_n)\log(1-\theta_0)$ is a convex function of $\theta_0$ (a sum of convex functions $-\log\theta_0$ and $-\log(1-\theta_0)$ with non-negative coefficients), so its sublevel set
$C_n$ is an interval.
\end{solution}

\section*{Sources and further reading}

\begin{itemize}
\item \emph{History.} J.~Ville, \emph{\'Etude critique de la notion de collectif}, Gauthier-Villars, 1939 (the word ``martingale'' and the inequality). L.~Bienvenu, G.~Shafer and A.~Shen, ``On the history of martingales in
the study of randomness'', \emph{Electronic J.\ History Probab.\ Statist.}\ 5(1) (2009), for Ville's thesis, its defence in March 1939 and his war. R.~Mansuy, ``The origins of the word `martingale' '', same issue
(the 1762 dictionary of the Acad\'emie fran\c caise is quoted there). J.~L.~Doob, \emph{Stochastic Processes}, Wiley, 1953. P.~L\'evy, \emph{Th\'eorie de l'addition des variables al\'eatoires}, Gauthier-Villars, 1937.
\item \emph{De Finetti.} B.~de Finetti, ``Funzione caratteristica di un fenomeno aleatorio'', \emph{Atti R.\ Accad.\ Naz.\ Lincei, Mem.}\ (6) 4 (1931), 251--299; ``La pr\'evision: ses lois logiques, ses sources
subjectives'', \emph{Ann.\ Inst.\ H.\ Poincar\'e} 7 (1937), 1--68; \emph{Theory of Probability}, Wiley, 1974, preface (``PROBABILITY DOES NOT EXIST''). E.~Hewitt and L.~J.~Savage, ``Symmetric measures on Cartesian
products'', \emph{Trans.\ Amer.\ Math.\ Soc.}\ 80 (1955), 470--501. D.~Aldous, ``Exchangeability and related topics'', \emph{\'Ecole d'\'et\'e de Saint-Flour XIII}, Lecture Notes in Math.\ 1117, Springer, 1985: the
general theory.
\item \emph{Bayesian consistency.} J.~L.~Doob, ``Application of the theory of martingales'', in \emph{Le Calcul des Probabilit\'es et ses Applications}, Colloques Internationaux du CNRS 13 (1949), 23--27.
D.~A.~Freedman, ``On the asymptotic behavior of Bayes' estimates in the discrete case'', \emph{Ann.\ Math.\ Statist.}\ 34 (1963), 1386--1403. W.~R.~Thompson, ``On the likelihood that one unknown probability
exceeds another in view of the evidence of two samples'', \emph{Biometrika} 25 (1933), 285--294.
\item \emph{Kakutani.} S.~Kakutani, ``On equivalence of infinite product measures'', \emph{Ann.\ of Math.}\ 49 (1948), 214--224.
\item \emph{Anytime-valid inference.} D.~A.~Darling and H.~Robbins, ``Confidence sequences for mean, variance, and median'', \emph{Proc.\ Nat.\ Acad.\ Sci.\ USA} 58 (1967), 66--68. H.~Robbins, ``Statistical methods related
to the law of the iterated logarithm'', \emph{Ann.\ Math.\ Statist.}\ 41 (1970), 1397--1409. S.~R.~Howard, A.~Ramdas, J.~McAuliffe and J.~Sekhon, ``Time-uniform, nonparametric, nonasymptotic confidence sequences'',
\emph{Ann.\ Statist.}\ 49 (2021), 1055--1080. A.~Ramdas, P.~Gr\"unwald, V.~Vovk and G.~Shafer, ``Game-theoretic statistics and safe anytime-valid inference'', \emph{Statist.\ Sci.}\ 38 (2023), 576--601.
\item D.~Williams, \emph{Probability with Martingales}, Cambridge University Press, 1991. The friendliest book on the subject; uniform integrability, backward martingales and Kakutani's theorem are all there.
\item G.~Shafer and V.~Vovk, \emph{Game-Theoretic Foundations for Probability and Finance}, Wiley, 2019. Probability rebuilt on Ville's betting idea.
\end{itemize}
